\chapter[Conclusão]{Conclusão}

 % Retomada do problema central
O presente trabalho abordou o problema relacionado à identificação e avaliação da dor em animais utilizados para pesquisas biomédicas, temática que permanece como um desafio relevante no contexto experimental devido às suas implicações diretas no bem-estar animal, na ética científica e na confiabilidade dos resultados obtidos. 

Neste contexto, foi desenvolvido um sistema computacional capaz de classificar automaticamente padrões faciais associados à presença e intensidade da dor utilizando a \acs{MGS} como referência. Paralelamente, foi feita uma investigação comparativa entre diferentes arquiteturas de aprendizado profundo, incluindo modelos convolucionais pertencentes às famílias \textit{ResNet}, \textit{DenseNet} e \textit{MobileNet}, além do modelo baseado em Vision Transformer, permitindo analisar como diferentes estratégias arquiteturais influenciam o desempenho da tarefa de classificação.

Dentre as arquiteturas avaliadas, os modelos convolucionais apresentaram desempenho superior ao \acs{VIT} nas condições experimentais consideradas. Esse resultado se justifica devido à demanda natural deste modelo por um grande volume de dados. Considerando a dimensão da base de dados utilizada nesta pesquisa, a arquitetura \acs{VIT} não conseguiu explorar plenamente seu potencial de generalização, comportamento que também foi evidenciado nos experimentos descritos na Seção 4.5.1, os quais demonstraram a influência da quantidade de amostras sobre o desempenho do modelo.

A análise comparativa entre as arquiteturas convolucionais evidenciou que os modelos baseados em estratégias de propagação e reaproveitamento de características, representados pelas famílias \textit{ResNet} e \textit{DenseNet}, apresentaram os resultados mais consistentes ao longo dos experimentos, conforme indicado na Tabela \ref{tab:comparacao_geral}. Tal comportamento sugere que mecanismos como conexões residuais e densas favorecem a preservação e o refinamento das representações aprendidas ao longo da rede, contribuindo para uma extração de características mais eficiente no problema investigado.

Considerando que a \acs{MGS} é fundamentada em um conjunto reduzido de unidades de ação faciais bem definidas e que as imagens utilizadas apresentam baixa complexidade visual, os resultados indicam que a preservação das informações extraídas nas camadas iniciais impactam  na identificação dos padrões discriminativos associados à presença e à intensidade da dor. Dessa forma, arquiteturas capazes de reutilizar essas representações ao longo do processo de aprendizado mostraram-se adequadas para o contexto experimental desta pesquisa.

Os resultados obtidos para a \textit{MobileNet}, por sua vez, também evidenciam que arquiteturas leves podem representar alternativas viáveis, mesmo quando não apresentam o melhor desempenho entre os modelos avaliados. A acurácia de aproximadamente 86\% obtida no conjunto de validação demonstra que sua capacidade de generalização permanece satisfatória, embora inferior à observada para os demais modelos convolucionais. O comportamento de \textit{overfitting} identificado ao longo do treinamento indica que existe margem para aprimoramento por meio da adoção de estratégias como \textit{transfer learning}, técnicas adicionais de regularização e ampliação da base de dados. Considerando seu reduzido custo computacional e sua facilidade de implantação em dispositivos com recursos limitados, a \textit{MobileNet} permanece como uma arquitetura promissora para aplicações práticas com possíveis limitações computacionais.

Em síntese, os resultados desta pesquisa demonstram a viabilidade no processo de automatização da avaliação da dor em camundongos. A utilização das técnicas apresentadas contribuem para o aumento a padronização das análises e reprodutibilidade dos estudos pré-clínicos, apoiando pesquisadores e médicos-veterinários durante o monitoramento dos modelos experimentais. Além de seus impactos sobre o bem-estar animal, essas contribuições favorecem a obtenção de dados experimentais mais consistentes para a pesquisa translacional  dos conhecimentos produzidos voltados à medicina humana.

\section{Principais Contribuições}

Os resultados obtidos ao longo deste estudo reforçam a viabilidade da aplicação de técnicas de aprendizado profundo e visão computacional para a avaliação automatizada da dor em camundongos, utilizando como referência os indicadores definidos pela \acs{MGS}. Nesse contexto, a principal contribuição desta pesquisa consiste na investigação de diferentes paradigmas arquiteturais aplicados ao problema de classificação da dor, permitindo identificar características, vantagens e limitações associadas a cada abordagem.

A análise comparativa conduzida entre arquiteturas convolucionais tradicionais e modelos baseados em mecanismos de autoatenção forneceu evidências importantes acerca da influência arquitetural sobre o desempenho da tarefa proposta. Os experimentos realizados permitiram avaliar como diferentes estratégias de extração, propagação e representação de características impactam o processo de aprendizado, contribuindo para uma melhor compreensão da adequação desses modelos a problemas de classificação biomédica baseados em imagens.

Entre as contribuições técnicas do trabalho, destaca-se o desenvolvimento de uma versão customizada da arquitetura \textit{DenseNet}, concebida com o objetivo de preservar os benefícios associados ao reaproveitamento denso de características, ao mesmo tempo em que reduz a complexidade computacional da configuração original \textit{DenseNet-121}. Os resultados obtidos demonstraram que essa adaptação foi capaz de manter elevado desempenho preditivo, constituindo uma alternativa potencialmente interessante para aplicações sujeitas a restrições computacionais.

Adicionalmente, os experimentos conduzidos com a arquitetura \acs{VIT} contribuíram para ampliar a compreensão sobre o comportamento de modelos baseados em \textit{Transformers} aplicados a cenários caracterizados por conjuntos de dados moderados e altamente especializados. Embora a arquitetura não tenha alcançado o melhor desempenho entre os modelos avaliados, os resultados obtidos fornecem subsídios importantes para futuras investigações envolvendo mecanismos de autoatenção aplicados à avaliação automatizada da dor.

Por fim, as contribuições deste trabalho extrapolam o contexto computacional, alcançando também aplicações relacionadas ao bem-estar animal e à pesquisa biomédica. A utilização de sistemas automatizados para avaliação da dor pode contribuir para a redução da subjetividade inerente às análises humanas, aumentar a reprodutibilidade experimental e auxiliar no monitoramento mais consistente das condições fisiológicas dos animais utilizados em pesquisa.


% SEÇÃO ALTERADA PARA COMENTAR SOBRE ELEMENTOS AUSENTES NO TRABALHO
\section{Discussões e Sugestões para Trabalhos Futuros}

Uma limitação importante desta pesquisa está relacionada ao tamanho e à diversidade da base de dados empregada nos experimentos. Embora o conjunto de imagens utilizado tenha se mostrado suficiente para investigar o comportamento das diferentes arquiteturas avaliadas e demonstrar a viabilidade da aplicação de técnicas de aprendizado profundo na avaliação automatizada da dor, sua dimensão permanece relativamente modesta quando comparada às bases de dados comumente utilizadas para o treinamento de modelos modernos de visão computacional. Nesse sentido, uma das principais perspectivas para trabalhos futuros consiste na ampliação da base de dados, tanto em quantidade de amostras, quanto na diversidade de informações, favorecendo o treinamento de modelos mais robustos e menos suscetíveis ao \textit{overfitting}. 

Adicionalmente, por mais que os experimentos conduzidos tenham permitido avaliar quantitativamente o desempenho das diferentes arquiteturas, não foi realizada uma análise das regiões da imagem que efetivamente contribuíram para as decisões dos classificadores. Nesse contexto, técnicas de Inteligência Artificial Explicável, como \textit{Grad-CAM}, \textit{Grad-CAM++} ou mecanismos equivalentes de visualização podem ser empregadas para gerar mapas de ativação (\textit{heatmaps}) capazes de evidenciar as regiões faciais consideradas mais relevantes durante o processo de classificação. Além de aumentar a interpretabilidade dos modelos, esse tipo de análise pode contribuir para verificar se as decisões tomadas pelas redes neurais estão efetivamente associadas às unidades de ação definidas pela \acs{MGS}.

Também é importante destacar que embora sejam utilizados conjuntos independentes de treinamento, validação e teste, o presente trabalho concentrou as análises sobre os conjuntos de treinamento e validação. Essa decisão foi motivada pelo objetivo central da pesquisa, que consistiu em investigar comparativamente o comportamento de diferentes arquiteturas de aprendizado profundo, e não na proposição de um modelo definitivo para implantação em ambientes reais. Entretanto, é inegável que a utilização de um conjunto independente de teste representa uma etapa importante para a validação final de modelos destinados à aplicação prática, de modo a apresentar uma oportunidade natural para a continuidade desta pesquisa.

Um outro aspecto metodológico refere-se à utilização do \textit{transfer learning}. Embora essa estratégia seja amplamente empregada na literatura para acelerar a convergência do treinamento e melhorar a capacidade de generalização dos modelos, optou-se por não empregá-la nos experimentos conduzidos com as arquiteturas convolucionais, buscando preservar as condições equivalentes de comparação entre os modelos investigados e minimizar a influência de representações previamente aprendidas em bases de dados externas. Contudo, para investigações futuras, recomenda-se a implementação desta técnica nos modelos avaliados, de modo a contribuir para a produção de um sistema mais otimizado e generalista, mesmo com limitações de dados.

Entretanto, essa estratégia não pôde ser mantida para a arquitetura \acs{VIT}. Diferentemente das redes neurais convolucionais, modelos baseados em mecanismos de autoatenção apresentam maior dependência de grandes volumes de dados para aprender representações visuais robustas, uma vez que não incorporam, de forma explícita, os vieses indutivos presentes nas operações convolucionais. Considerando as limitações do conjunto de dados disponível neste trabalho, a utilização de pesos pré-treinados tornou-se necessária para viabilizar o treinamento da arquitetura, permitindo explorar seu potencial de representação mesmo em um cenário de dados mais restrito.

Por fim, a continuidade natural deste trabalho também envolve o desenvolvimento de sistemas de inferência em tempo real voltados para aplicações práticas em ambientes laboratoriais. A integração dos modelos treinados a dispositivos embarcados ou aplicações móveis pode viabilizar soluções automatizadas para monitoramento contínuo da dor em animais, reduzindo a dependência humana e contribuindo para práticas experimentais mais éticas, rápidas e padronizadas.

\section{Contribuições em Produção Bibliográfica}

%A produção bibliográfica associada ao presente trabalho foi impactada por fatores relacionados à reformulação do escopo original do projeto, aos trâmites de aprovação junto ao Comitê de Ética (\acs{CEUA}) e à dependência de equipes especializadas para a aquisição e validação da base de dados utilizada nos experimentos. Inicialmente, a pesquisa tinha como foco a previsão do sucesso de eclosão de ovos de galinha por meio da análise conjunta de imagens digitais e metadados associados. Entretanto, os resultados obtidos não demonstraram evidências suficientes para sustentar a continuidade dessa linha de investigação, motivando a redefinição do tema de pesquisa.

%A transição para o estudo da avaliação automatizada de dor em camundongos introduziu novos desafios metodológicos e operacionais. Além da necessidade de aprovação ética para utilização dos animais, a coleta e rotulação da base de dados dependeram da atuação de profissionais especializados da área veterinária, responsáveis tanto pelo manejo dos animais quanto pela identificação dos indicadores faciais associados à dor. Embora essenciais para garantir a qualidade científica e o cumprimento das normas de bem-estar animal, essas etapas ampliaram significativamente o tempo necessário para a execução da pesquisa.

A produção bibliográfica associada ao presente trabalho foi impactada por fatores relacionados à reformulação do escopo original do projeto, aos trâmites de aprovação junto ao Comitê de Ética (\acs{CEUA}) e à dependência de equipes especializadas para a aquisição e validação da base de dados utilizada nos experimentos. Apesar de tais limitações, os resultados obtidos culminaram na elaboração e submissão de um artigo científico ao Simpósio Brasileiro de Computação Gráfica e Processamento de Imagens (\acs{SIBGRAPI}), um dos principais eventos nacionais nas áreas de Visão Computacional, Processamento de Imagens e Computação Gráfica. O trabalho apresenta uma investigação específica sobre a aplicação de arquiteturas convolucionais para a detecção e quantificação automática da dor em camundongos com base nos indicadores definidos pela \textit{Mouse Grimace Scale}, incluindo uma análise comparativa entre diferentes modelos adotados.

Além da produção já submetida, os resultados alcançados ao longo desta pesquisa evidenciaram oportunidades para novas publicações científicas. Entre elas, destaca-se a ferramenta computacional desenvolvida para extração, organização e padronização da bases de imagens, concebida para auxiliar pesquisadores na construção de conjuntos adequados ao treinamento de modelos classificadores. Em razão de sua aplicabilidade potencial em diferentes domínios da visão computacional, encontra-se prevista a elaboração de um artigo específico abordando sua arquitetura, funcionalidades e aplicações.

Adicionalmente, os experimentos realizados abriram novas perspectivas de investigação relacionadas ao emprego de arquiteturas avançadas para análise automatizada da dor em animais. Nesse contexto, prevê-se a elaboração de uma publicação adicional explorando aspectos complementares aos resultados apresentados nesta dissertação, ampliando as discussões sobre metodologias, arquiteturas e aplicações futuras da abordagem proposta. Dessa forma, espera-se que os resultados obtidos nesta pesquisa contribuam não apenas para a conclusão do presente trabalho, mas também para a continuidade das atividades científicas, fomentando novas publicações e ampliando o impacto acadêmico das contribuições desenvolvidas.